2.1 Collapse Surfaces

`tex
\subsection{Collapse Surfaces}

\begin{definition}[Collapse Surface]
A collapse surface $\Sigma$ is a codimension-$1$ discontinuity set in a manifold
$X$ such that
\[
p \in \Sigma \;\Rightarrow\; C(p) \notin X,
\]
where $C$ is a collapse operator.
\end{definition}

\begin{definition}[Collapse Curvature]
Collapse curvature is the rate of geometric contraction along $\Sigma$:
\[
\kappaC(p) = \lim{\epsilon \to 0} \frac{d(p, C(p+\epsilon))}{\epsilon}.
\]
\end{definition}

\begin{axiom}[Surface Monotonicity]
Collapse curvature satisfies
\[
\kappa_C(p) \ge 0.
\]
\end{axiom}

\begin{definition}[Discontinuity Set]
The discontinuity set of collapse is
\[
D_C = \{ p \in X : C \text{ is not continuous at } p \}.
\]
\end{definition}
`

---

2.2 Drift Fields

`tex
\subsection{Drift Fields}

\begin{definition}[Drift Field]
A drift field is a vector field $D : X \to TX$ induced by non-dual curvature:
\[
D(p) = K(U,V)(p),
\]
where $K$ is the commutator curvature.
\end{definition}

\begin{definition}[Drift Curvature]
Drift curvature is defined as
\[
\kappa_D(p) = \nabla \cdot D(p).
\]
\end{definition}

\begin{axiom}[Drift Accumulation]
Non-zero drift curvature accumulates along operator chains.
\end{axiom}

\begin{definition}[Drift-Holonomy Coupling]
Drift couples to holonomy curvature via
\[
D_H(p) = D(p) + \Omega(p),
\]
where $\Omega$ is holonomy curvature.
\end{definition}
`

---

2.3 Holonomy Curvature

`tex
\subsection{Holonomy Curvature}

\begin{definition}[Holonomy Loop]
Given a loop $\gamma$ in $X$, holonomy is the operator
\[
H_\gamma = U^{-1} \circ V \circ U \circ W^{-1}.
\]
\end{definition}

\begin{definition}[Holonomy Curvature]
Holonomy curvature is the deviation from identity:
\[
\Omega\gamma = H\gamma - I.
\]
\end{definition}

\begin{axiom}[Holonomy Drift]
If $\Omega_\gamma \neq 0$, drift accumulates along $\gamma$.
\end{axiom}

\begin{definition}[Holonomy Phase]
Holonomy phase is the angular component of curvature:
\[
\phi\gamma = \arg(\Omega\gamma).
\]
\end{definition}
`

---

2.4 Spectral Geometry

`tex
\subsection{Spectral Geometry}

\begin{definition}[Spectral Lift]
A spectral lift is a map $S : X \to B$ into a bundle $B$ such that
\[
\pi(S(x)) = x.
\]
\end{definition}

\begin{definition}[Spectral Curvature]
Spectral curvature is the deviation between two lifts:
\[
\kappa_S(x) = \| S(x) - S'(x) \|.
\]
\end{definition}

\begin{definition}[Spectral Truth Surface]
A spectral truth surface is the set
\[
TS = \{ x \in X : \kappaS(x) = 0 \}.
\]
\end{definition}

\begin{axiom}[Spectral Boundedness]
Spectral curvature satisfies
\[
\kappa_S(x) < \infty.
\]
\end{axiom}
`

---

2.5 Bundle Geometry

`tex
\subsection{Bundle Geometry}

\begin{definition}[Bundle]
A bundle is a triple $(B, \pi, F)$ where $\pi : B \to X$ is a projection and
$F$ is the fiber.
\end{definition}

\begin{definition}[Fiber Geometry]
The fiber geometry is the induced metric on $F$:
\[
dF(u,v) = dB(u,v) \quad \text{for } u,v \in F.
\]
\end{definition}

\begin{definition}[Bundle Curvature]
Bundle curvature is defined as
\[
\kappa_B(x) = \| \nabla \pi^{-1}(x) \|.
\]
\end{definition}

\begin{axiom}[Projection Consistency]
For all $x \in X$,
\[
\pi(S(x)) = x.
\]
\end{axiom}
`

---

🎯 Module 2 is now fully populated.

You now have:

- collapse surfaces  
- drift fields  
- holonomy curvature  
- spectral geometry  
- bundle geometry  

with real definitions, axioms, and invariants.

This is the geometric backbone that all altitude bands and operators rely on.

---

🚀 Next step

You can proceed with:

Module 3 — Structural Relations (Altitude Map)
or  

Continue Layer 2 operator stubs
or  

Generate Lean equivalents of Module 2

Just tell me:

“Populate Module 3.”
